\documentclass[10pt,a4paper]{amsart}
\usepackage[margin=19mm]{geometry}
\usepackage{amsmath,amssymb,mathtools}
\usepackage[scr=rsfs,cal=euler]{mathalfa}
\usepackage{xfrac}
\usepackage[unicode,hidelinks,bookmarks=false]{hyperref}
\newcommand{\ie}{i.e.\ }
\newcommand{\cf}{cf.\ }
\newcommand{\resp}{respectively\ }
\newcommand{\Subsection}{\subsection}
\providecommand{\nobreakdash}{\nobreak\hbox{-}}
\setlength{\emergencystretch}{2em}
\multlinegap0pt
\usepackage{bm}
\usepackage{bbm}
\usepackage{dsfont}
\usepackage{xspace}
\usepackage{cancel}
\usepackage{mathtools}
\usepackage{amsthm}
\makeatletter
\@ifundefined{proposition}{\newtheorem{proposition}{Proposition}}{}
\@ifundefined{lemma}{\newtheorem{lemma}[proposition]{Lemma}}{}
\makeatother
\newcommand{\ZZ}{\mathbf{Z}}                          % /Z
\title[Normal Subgroup Theorem for groups acting on $\tilde A\_2$-bu]{Normal Subgroup Theorem for groups acting on $\tilde A\_2$-buildings}
\author{Uri Bader, Alex Furman, Jean L\'ecureux}
\date{Source: arXiv:2310.13472v2 (CC BY 4.0)}
\begin{document}
\maketitle

\noindent\emph{Source and licence.} Normal Subgroup Theorem for groups acting on $\tilde A_2$-buildings. Version arXiv:2310.13472v2 (CC BY 4.0), distributed under the Creative Commons Attribution 4.0 International licence, as recorded in the arXiv licence metadata for this exact version and shown on the abstract page https://arxiv.org/abs/2310.13472v2. Full rights notice: licenses/arxiv-2310.13472-CC-BY-4.0.txt.\par\smallskip

\noindent\textit{Editorial scope.} Counted unit: the Lemma whose source label is \texttt{lem:seqextensions} in the article (source0.tex lines 980--982, source label \texttt{lem:seqextensions}), delivered together with the article's own complete original proof of it (source0.tex lines 984--1012, one \texttt{proof} environment of 29 lines). No mathematics is written, repaired or re-derived: statement and proof are the article's own text line for line. Required context retained uncounted: lem:extstd (lines 940--944). Every definition, display and label of those lines is retained unchanged. Omitted from this package and not redistributed: 1--939, 945--979, 983--983, 1013--2261. Nothing inside a retained excerpt is omitted or altered: every retained body is delivered exactly as the article's lines are written, so no omission receipt of any kind accompanies this package. Citations: the retained excerpts carry no citation command at all, so no bibliography entry of the article is transcribed and no reference is redistributed. Reference adaptation: none was needed and none was made; every reference inside the delivered unit resolves inside it, so no unresolved reference or citation is produced. Printed slips kept literal: no printed word, symbol, hypothesis or number of the counted statement or of its proof was altered. Layout and class: the article's own class is \texttt{article} with packages including fontenc, inputenc, url, hyperref, amsfonts, amsmath, amssymb, amsthm, enumerate, amsrefs, comment, array, imakeidx, mathrsfs; the wrapper is the lane's pinned \texttt{amsart} editorial wrapper at 10pt on A4, which restates the theorem-like environments the retained text needs and adds the article's own macro definitions that the retained text needs (\texttt{\textbackslash ZZ}), with extra wrapper packages (bm, bbm, dsfont, xspace, cancel, mathtools). The delivered numbering is the unit's own. Assets: no retained span contains a figure environment, a graphics-inclusion command, a PDF-page inclusion, a table or a TikZ picture; no image file is copied into this package, the package declares no asset dependency and no image file is redistributed with it. Licence: version arXiv:2310.13472v2 of this article is distributed under the Creative Commons Attribution 4.0 International licence, as recorded in the arXiv licence metadata for this exact version and shown on the abstract page https://arxiv.org/abs/2310.13472v2. A full rights notice is delivered at licenses/arxiv-2310.13472-CC-BY-4.0.txt. Characters: every retained excerpt is pure ASCII, so the delivered engine, XeLaTeX with its default Latin Modern text font, reports no missing character.
\begin{lemma}\label{lem:extstd}
Let $Y\subset Y'$ be an isometric inclusion of crazy diamonds. Let $Y'=Y(T',x',y',s',t')$. Then there exists a subtree $T\subset T'$, $x,y\in T$ and $s,t\in \ZZ$ such that $Y=Y(T,x,y,s,t)$ and the inclusion $Y\subset Y'$ is as described above.

%Furthermore, $T,x,y,s,t$ are uniquely defined by the inclusion $X\subset X'$. 
\end{lemma}

\begin{lemma}\label{lem:seqextensions}
Let $Y\subset Y'$ be an isometric inclusion of crazy diamonds, and assume that $Y$ is not degenerate. Then there is a finite sequence $Y=Y_0,Y_1,\dots,Y_n=Y'$ of crazy diamonds such that $Y_i\subset Y_{i+1}$ is an elementary extension. %Furthermore, the number $n$ (and in fact the number of extensions of each type) is uniquely defined.
\end{lemma}

\begin{proof}

As in Lemma \ref{lem:extstd}, we can and shall assume that $T\subset T'$. Note that if $z\in T'$ has valency $>1$ then $l'(z')\geq 1$.
Also note that since $c'(x)\geq c(x)$, we have $s'\geq s+d(x,x')\geq s$. 


%Since extensions of type $C$ are the only ones which can (and do) increase $s$, 
%each of them by 1, if there is a sequence of elementary extensions from $Y$ to $Y'$, then it contains exactly $s'-s$  extension of type $C$. 
%Similarly a sequence of elementary extensions from $Y$ to $X'$ will contain $t-t'$ extensions of type $F$.
%The extensions of type $B_z$, for $z\in T'$ are the only ones that can change the tree $T$, and only by adding an edge. Therefore we need to perform such an extension exactly the number of edges in $T'\setminus T$. More precisely, for each $z\in T'$, we need to do $v'(z)-v(z)$ extensions of type $B_z$, where $v'(z)$ (resp. $v(z)$) is the valency of $z$ in $T'$ (resp. in $T$ if $z\in T$, and $1$ otherwise). 

%The only thing to do is to check that there exists indeed a sequence of elementary extensions from $Y$ to $Y'$.

First assume that $Y$ is non-degenerate.
 Our sequence of extensions will proceed as follows. The notations relative to $Y_i$ will be the same as for $Y$ but with index $i$. We use $i$ as a temporary variable and change it after each step.

 By doing two successive extensions of type $C$, we can replace $s$ by $s+2$ without changing $x$. We start by this move, doing $d$ extensions of type $C$, in order to get that $s_i=c_i(x)=c'(x)$. Note that $c'(x)-c(x)$ must be even since $Y$ is a subcomplex of $Y'$.  
  Then we do in a similar way the right number of extensions of type $F$, in order to get that $t_{i}=f_{i}(y)=f'(y)$. It follows that $c_i(z)=c'(z)$ as soon as $d(z,x)<d(z,x')$. By doing an extension of type $C$ again, moving $x$ to the vertex $z_0$ adjacent to $x$ and closer to $x'$, we get that  $c_i(z)=c'(z)$ for $z$ satisfying $d(z,x)<d(z,x')+1$. Repeating the same thing as many times as necessary, we can get that $c_i(z)=c'(z)$ for every $z\in T$. Similarly we can obtain $f_i(z)=f'(z)$ for every $z\in T$. 
  
  If $z\in T$ is of valency $>1$ in $T'$, then $l_i(z)=l'(z)\geq 1$. Therefore we can perform an extension of type $B_z$, in order to add as many edges as necessary. Repeating the above steps, we see that we can get that $T_i=T'$, $c_i=c'$ and $f_i=f'$, which concludes the argument.

Then let us assume that $Y$ is vertical (but not a point). If $Y'$ is non-degenerate, then there is an alcove in $Y'$ which is adjacent to some panel of $Y$; let $Y_1$ be the convex hull of $Y'$ and this alcove. Then $Y\subset Y_1$ is an elementary extension of type $B$, and by the previous argument there exists a finite sequence of elementary extensions from $Y_1$ to $Y'$.
If $Y'$ is also degenerate, then it must be vertical, then it is just a segment containing $Y$, and one can pass from $Y$ to $Y'$ by adding edges to this segment, that is, by a finite sequence of elementary extensions of type $C_v$ or $F_v$.

Now let us assume that $Y$ is horizontal (and not a point). If $Y'$ is non-degenerate, again there is an alcove in $Y'$ which is adjacent to $Y$, and the convex hull $Y_1$ of this alcove and $Y$ is an elementary extension of $Y$ of type either $C$ or $F$.  If $Y'$ is degenerate, then it must be horizontal, and then again one can pass from $Y'$ to $Y$ by a sequence of elementary extensions of type $B^f$ or $B^c$.

Finally, if $Y$ is a point, then there exists a segment in $Y'$ containing $Y$, and  the inclusion of $Y$ in this segment is
an elementary extension of type either $C_v$, $F_v$, $B^f$ or $B^c$ depending on the position of this segment in $Y'$. In any case applying the previous argument we can construct a sequence of elementary extensions from $Y$ to $Y'$.
\end{proof}

\end{document}
